\chapter{Anal Fissure}

\section*{Definition}
An anal fissure is a small tear in the skin lining the anal canal, usually causing sharp pain during and after a bowel movement and bright-red bleeding. Most occur in the posterior midline; anterior fissures are more common after childbirth. An acute fissure is recent, while a chronic fissure persists or recurs despite initial care, often for more than 6--8 weeks.

\section*{When to See a Doctor}

\begin{itemize}
 \item Severe anal pain during and after defecation lasting minutes to hours, especially if the diagnosis is uncertain
 \item Bright red blood on toilet paper or coating the stool surface
 \item Fissure persisting beyond 6--8 weeks, recurring, or not improving with stool-softening measures
 \item Multiple fissures or fissures in atypical locations (lateral, anterior) suggesting underlying disease
 \item Fever, pus, a painful lump, heavy bleeding, weight loss, abdominal pain, diarrhea, or a change in bowel habits
\end{itemize}

\section*{How It Develops}
Passage of hard stools, straining, or frequent diarrhea can injure the anal lining. Pain may trigger increased internal sphincter tone, which can reduce local blood flow and delay healing. Some fissures have different causes, especially when they are multiple, lateral, recurrent, or associated with another illness.

\section*{Causes}
\begin{itemize}
 \item \textbf{Constipation:} Passage of hard stools, straining during defecation (most common cause)
 \item \textbf{Chronic diarrhea:} Repetitive trauma from frequent liquid stools
 \item \textbf{Childbirth:} Vaginal delivery causing perineal trauma and anal sphincter injury
 \item \textbf{Inflammatory bowel disease:} Crohn's disease (especially atypical lateral or multiple fissures), ulcerative colitis
 \item \textbf{Other:} Anal intercourse, colonoscopy trauma, tuberculosis, syphilis, HIV, leukemia, sexual abuse
\end{itemize}

\section*{Related Symptoms}
\begin{itemize}
 \item Sharp anal pain during defecation, post-defecatory throbbing pain, bright red blood per rectum, itching ani
\end{itemize}
\textbf{Important:} A sentinel skin tag or enlarged anal papilla can occur with chronic fissures, but their presence does not prove that healing without surgery is impossible. Atypical fissures need evaluation for inflammatory, infectious, traumatic, or malignant causes.

\section*{Medical Evaluation}
\begin{itemize}
 \item Visual inspection with gentle separation of the buttocks is often sufficient initially. A digital examination or anoscopy may be deferred if pain is severe and should be performed gently when needed.
 \item Identify fissure location (posterior midline most common), depth, skin tag, hypertrophied papilla
 \item Anoscopy or proctoscopy for fissures in atypical locations to exclude Crohn disease, malignancy, or infection
 \item Consider evaluation for inflammatory bowel disease or infection when there are lateral or multiple fissures, fistulas, diarrhea, abdominal symptoms, or systemic features; colonoscopy is selected according to the presentation
 \item Biopsy of non-healing or suspicious fissure edges in immunocompromised or at-risk people
\end{itemize}

\section*{Treatment Options}
\begin{itemize}
 \item Conservative care with fiber, appropriate fluids, stool-softening treatment, and warm baths is used for acute fissures and may be continued for chronic disease.
 \item Topical calcium-channel blockers such as diltiazem or nifedipine can relax the internal sphincter and generally cause fewer headaches than topical nitrates.
 \item Topical nitroglycerin is another option, but headaches and low blood pressure can limit its use; it must not be combined with phosphodiesterase-5 inhibitors.
 \item Botulinum toxin injection may be considered for chronic fissures and has similar effectiveness to topical therapy in many studies.
 \item Tailored lateral internal sphincterotomy is highly effective for selected chronic fissures but carries a risk of fecal incontinence; sphincter-preserving options may be preferred for people at higher risk.
\end{itemize}

\section*{Self-Care and Home Management}
\begin{itemize}
 \item Sitz baths (warm water, no soap) for 10--15 minutes, 2--4 times daily and after each bowel movement
 \item Use fiber or a clinician-recommended laxative such as psyllium or polyethylene glycol to keep stools soft; the best choice depends on constipation, diarrhea, and other conditions
 \item Increase dietary fiber gradually and drink enough fluid to avoid dehydration, unless a clinician has prescribed fluid restriction
 \item Topical lidocaine 2--5\% gel or ointment applied before bowel movements for pain control
 \item Avoid straining during defecation and prolonged sitting on the toilet
\end{itemize}

\section*{References}
\begin{thebibliography}{99}
\bibitem{anal_fissure:ref1} Beaty JS, Shashidharan M. ``Anal fissure.'' \textit{Clin Colon Rectal Surg}. 2016;29(1):30-37.
\bibitem{anal_fissure:ref2} Wald A. ``Diagnosis and management of anal fissure.'' \textit{Am Fam Physician}. 2017;95(4):234-240.
\bibitem{anal_fissure:ref3} Nelson RL, Thomas K, et al. ``Nonsurgical therapy for anal fissure.'' \textit{Cochrane Database Syst Rev}. 2012;(2):CD003431.
\bibitem{anal_fissure:ref4} Beck DE, Roberts PL, et al. \textit{The ASCRS Textbook of Colon and Rectal Surgery}, 3rd ed. Springer, 2018.
\bibitem{anal_fissure:ref5} UpToDate. ``Anal fissure: medical management.'' Wolters Kluwer, 2023.
\bibitem{anal_fissure:ref6} Davids JS, et al. The American Society of Colon and Rectal Surgeons clinical practice guidelines for the management of anal fissures. \textit{Dis Colon Rectum}. 2023;66:190--199.
\end{thebibliography}
